\section{Introduction}
\label{sec:introduction}

This report begins from a small but persistent bottleneck in synthetic data construction. To build scene-aware object rotation pairs, each object must be rendered in a fixed environment with stable lighting, grounding, scale, viewpoint, and background consistency. In practice, this requires repeated Blender-based tuning \cite{blenderFoundationBlender}: render samples, inspect failures, diagnose whether the problem comes from placement, lighting, camera framing, material quality, or reconstruction artifacts, adjust parameters or scripts, and render again. The bottleneck is not merely that Blender is inconvenient. The deeper issue is that manual synthetic-data construction is already a human-in-the-loop closed process, but the loop is often implicit, slow, and difficult to scale.

This pain point becomes costly when the goal is modest but task-specific. A user may only want to add a small batch of new objects or quickly build a dataset for a new object-level edit, yet still needs rendering expertise, repeated trial and error, and sometimes exploratory code-level tuning through Blender MCP-style tooling or direct code edits. Without an explicit workflow, the same failures recur across objects and dataset rounds. Worse, if downstream training regresses, it is hard to trace whether the cause was a bad asset, a poor render, a weak instruction, a split issue, or the training configuration.

The central observation is that synthetic dataset construction should not be modeled as a one-shot pipeline. It is an iterative workflow: observe, diagnose, adjust, regenerate, and validate. This report proposes a \method{} that makes this loop explicit. The engine turns a dataset request into staged artifacts, review signals, feedback actions, and acceptance outcomes. VLM/CV review and AI coding agents replace part of the human-driven Blender tuning process, while downstream training and evaluation provide feedback at the dataset-round level.

The framework is evaluated on the original motivating task: scene-aware object rotation data construction. This task is appropriate because it requires stable background context, object identity preservation, viewpoint control, and lighting and grounding consistency. It also forces a clear distinction between object rotation and camera-orbit multiview rendering. The case study is therefore not the general framework itself, but a concrete validation setting for testing whether the workflow can produce useful train-ready data and diagnose its own bottlenecks.

The main finding is diagnostic: the workflow surfaces useful weak-angle signal while exposing an unresolved render-quality bottleneck, demonstrating that both loops actively shape the outcome. Concrete metrics are deferred to \cref{sec:results}.

This report makes three contributions:
\begin{itemize}[leftmargin=1.5em]
  \item It reframes autonomous visual dataset construction as a closed-loop workflow-as-skill data engine rather than a one-shot generation pipeline.
  \item It defines a dual-loop design in which VLM/CV feedback and AI-agent actions support generation-time self-correction, while downstream probes support validation-time self-improvement.
  \item It instantiates the framework on scene-aware object rotation editing and shows that a mixed first feedback-scaling round can become actionable diagnostic evidence rather than wasted trial and error.
\end{itemize}
